\subsection{The full Dedekind gate as a Milnor residue condition}

\begin{lemma}[Proved; the top-degree residue obstruction]
\label{A:Dedekind-Milnor-obstruction}
Let $R$ be a Dedekind domain, let $p\in R^\times$ be prime, and
put $K=\Frac R$ and $n=\cd_p(K)<\infty$. Suppose that $R$ is not a
field, so $n\geq1$. For each finite separable extension $L/K$, let
$R_L$ be the integral closure of $R$ in $L$, and put $Y_L=\Spec R_L$.
Then $R_L$ is a finite $R$-algebra and a Dedekind domain. Write
$Y_L^{(1)}$ for its closed points and $\partial_y$ for the tame
residue in Milnor $K$-theory at the discrete valuation of $y$.
There is a canonical isomorphism
\[
 H^{n+1}(Y_L,\mu_p^{\otimes n})\simeq
 \operatorname{coker}\left[
 K_n^M(L)/p\xrightarrow{(\partial_y)_y}
 \bigoplus_{y\in Y_L^{(1)}}K_{n-1}^M(\kappa(y))/p
 \right].
\]
Here $K_0^M(E)=\Z$. No excellence hypothesis is imposed.
\end{lemma}
\begin{proof}
We first justify the normalization assertion rather than assume a
Japanese hypothesis. Choose a $K$-basis $b_1,\ldots,b_d$ of $L$
consisting of elements integral over $R$, by clearing the finitely
many denominators in the equations of an arbitrary basis. If
$z\in R_L$, every $\operatorname{Tr}_{L/K}(zb_i)$ is integral over
$R$ and belongs to $K$, hence belongs to $R$ because $R$ is normal.
The nondegenerate trace pairing for the separable extension $L/K$
places $R_L$ inside the finite free $R$-lattice dual to
$\sum_i Rb_i$. A submodule of a finite module over the noetherian
ring $R$ is finite. Thus $R_L$ is finite; it is normal and
one-dimensional, so is Dedekind.

Finite separable extension gives $\cd_p(L)\leq n$. For every closed
point $y$ the discrete-valuation lower bound gives
$\cd_p\kappa(y)\leq n-1$. When $p=2$, the residue field cannot be
orderable: an ordering would extend through the regular local
ring $(R_L)_y$ to $L$, contrary to $\cd_2(L)<\infty$; this is the
regular-local case proved in Proposition~\ref{A:realexception}.

For a finite set $Z\subset Y_L^{(1)}$, localization and codimension-one
purity give the residue exact sequence. Taking the filtered colimit
over $Z$, using $\bigcap_Z(Y_L-Z)=\Spec L$ and cohomological
continuity, gives
\[
 H^n(L,\mu_p^{\otimes n})\longrightarrow
 \bigoplus_{y\in Y_L^{(1)}}
 H^{n-1}(\kappa(y),\mu_p^{\otimes(n-1)})
 \longrightarrow H^{n+1}(Y_L,\mu_p^{\otimes n})
 \longrightarrow H^{n+1}(L,\mu_p^{\otimes n})=0.
\]
The purity used here is only the discrete-valuation case. On a
strict henselization the tame inertia calculation gives
$i^!\mu_p^{\otimes n}=\mu_p^{\otimes(n-1)}[-2]$; it requires
$p$ invertible, not excellence. It is the same local calculation
underlying \citet[Expos\'e X, Th\'eor\`eme 2.2]{A_SGA4}.

The norm-residue isomorphism theorem
\citep[Theorem 6.16, p.~436]{A_Voevodsky2011} identifies the first
two groups with the displayed Milnor groups. It respects the
residue arrows: the residue of a cup product of unramified unit
classes is zero, and the residue of
$(\pi)\cup(u_1)\cup\cdots\cup(u_{n-1})$ is
$(\overline u_1)\cup\cdots\cup(\overline u_{n-1})$, for a
uniformizer $\pi$ and units $u_i$. These are precisely the defining
formulas for the tame Milnor residue; multilinearity and the
Steinberg relation reduce every symbol to these cases. The
exact sequence therefore has the stated cokernel.
\end{proof}

\begin{theorem}[Proved; an exact criterion in every generic dimension]
\label{A:Dedekind-Milnor-criterion}
With the notation of Lemma~\ref{A:Dedekind-Milnor-obstruction},
the following are equivalent:
\begin{enumerate}
\item $\cd_p(\Spec R)\leq n$.
\item For every finite separable $L/K$ containing a primitive
$p$th root of unity, the residue homomorphism
\[
 K_n^M(L)/p\longrightarrow
 \bigoplus_{y\in(\Spec R_L)^{(1)}}K_{n-1}^M(\kappa(y))/p
\]
is surjective.
\end{enumerate}
Consequently A.2 for this Dedekind domain holds if and only if there
is a nonzero $a\in R$ for which condition (2) holds simultaneously
with $R_L$ replaced by $R_L[1/a]$. For $n=1$, condition (2) is
$\operatorname{Pic}(R_L)/p=0$. This extends the finite-separable
version of Corollary~\ref{A:Dedekind-Pic-criterion} without
excellence.
\end{theorem}
\begin{proof}
Suppose (1). For the finite morphism $f:Y_L\to\Spec R$, exact
finite pushforward gives
$H^q(Y_L,\mathcal E)=H^q(\Spec R,f_*\mathcal E)$ for every torsion
sheaf $\mathcal E$. Hence $\cd_p(Y_L)\leq n$, and the preceding
lemma proves (2).

Conversely, all residue fields of $R$ are nonorderable when $p=2$,
by the regular-local ordering argument just used. The preliminary
bound of Corollary~\ref{A:generic-plus-dimension} gives
$\cd_p(R)\leq n+1$. It is enough to prove that
$H^{n+1}(R,\mathcal F)=0$ for each constructible
$\mathbf F_p$-sheaf $\mathcal F$.

Choose a finite Galois extension $L/K$ that trivializes its generic
Galois module and contains $\mu_p$. Shrink to a nonempty open
$j:V\hookrightarrow\Spec R$ on which $\mathcal F$ is locally
constant and $f:Y_L\to\Spec R$ is \'etale. Put
$V_L=f^{-1}V$, let $j':V_L\hookrightarrow Y_L$, and write
$f_V:V_L\to V$ for the restricted finite \'etale map. A finite
free module over the group algebra of $\operatorname{Gal}(L/K)$
surjects onto the finite generic module. Thus, for some finite $r$,
it gives a morphism
\[
 \mathcal P=j_!(f_V)_*\mathbf F_p^r
       =f_*j'_!\mathbf F_p^r\longrightarrow\mathcal F
\]
that is surjective on $V$. Let $\mathcal N$ be its kernel,
$\mathcal I$ its image, and $\mathcal Q$ its cokernel.
The latter is supported on finitely many closed points and has
cohomology zero above $n-1$. The two short exact sequences with
middle terms $\mathcal P$ and $\mathcal F$, together with
$H^{n+2}(R,\mathcal N)=0$, show that
$H^{n+1}(R,\mathcal F)$ is a quotient of
$H^{n+1}(R,\mathcal P)$.

The finite-support complement $Y_L-V_L$ also has cohomological
dimension at most $n-1$. Therefore extension by zero and finite
pushforward give
\[
 H^{n+1}(R,\mathcal P)
 =H^{n+1}(Y_L,j'_!\mathbf F_p^r)
 =H^{n+1}(Y_L,\mathbf F_p)^r.
\]
Since $\mu_p\subset L$ and $R_L$ is integrally closed with $p$
invertible, $\mu_p$ is a constant sheaf on $Y_L$. Condition (2)
and Lemma~\ref{A:Dedekind-Milnor-obstruction} make the last group
zero. This proves (1), first for constructible $\mathbf F_p$
coefficients and then for every $p$-primary torsion sheaf by
Lemma~\ref{A:basic}.

On a smaller open, the direct sum in (2) is obtained by omitting
the removed primes; its surjectivity follows from that on the
larger open. If one $a$ as stated exists, the upper bound holds on
every smaller affine open; Proposition~\ref{A:spread} provides the
simultaneous lower bound. Conversely, A.2 gives an affine open
with the upper bound, and a principal nonempty open inside it.
Finite normalization commutes with this localization, so (2)
holds for $R_L[1/a]$. Finally, when $n=1$, the residue map is the
valuation divisor map $L^\times/(L^\times)^p\to
\bigoplus_y\Z/p$; its cokernel is exactly
$\operatorname{Pic}(R_L)/p$.
\end{proof}

\begin{corollary}[Proved; all semilocal Dedekind domains]
\label{A:semilocal-Dedekind}
Let $R$ be a semilocal Dedekind domain, let $p$ be invertible in
$R$, and suppose $\cd_p(\Frac R)<\infty$. Then
\[
 \cd_p(\Spec R)\leq\cd_p(\Frac R),
\]
and A.2 holds. Neither henselianity nor excellence is assumed.
\end{corollary}
\begin{proof}
Every finite separable normalization $R_L$ is again semilocal
Dedekind: its maximal ideals lie over the finitely many maximal
ideals of $R$, with finite fibres. Let them be
$\mathfrak m_1,\ldots,\mathfrak m_s$. By the Chinese remainder
theorem choose $\pi_i\in R_L$ with valuation one at
$\mathfrak m_i$ and valuation zero at every other maximal ideal,
prescribing its class modulo $\mathfrak m_i^2$ and modulo the
other maximal ideals. For a symbol
$\{\overline u_1,\ldots,\overline u_{n-1}\}$ in the residue field
at $\mathfrak m_i$, lift each $\overline u_j$ to an element
$u_j\in R_L$ with value one at all other maximal ideals. These
lifts avoid every maximal ideal, and are therefore global units.
The symbol $\{\pi_i,u_1,\ldots,u_{n-1}\}\in K_n^M(L)/p$ has the
desired residue at $\mathfrak m_i$ and zero residue everywhere
else. For $n=1$ this simply says that $\pi_i$ has valuation vector
the $i$th basis vector. Symbols generate the Milnor groups,
proving surjectivity for every $L$. The preceding theorem proves
the upper bound and A.2. If $R$ is a field, both claims are immediate.
\end{proof}

\begin{theorem}[Proved; normalization removes the nonnormal curve gate]
\label{A:curve-normalization}
Let $R$ be any one-dimensional noetherian domain, let $p$ be
invertible in $R$, let $K=\Frac R$ with $n=\cd_p(K)<\infty$,
and let $B$ be the normalization of $R$ in $K$. Then A.2 holds
for $\Spec R$ if and only if it holds for $\Spec B$.
Equivalently, it holds if and only if there is a nonzero $a\in R$
such that the residue maps of
Theorem~\ref{A:Dedekind-Milnor-criterion}, for the finite separable
normalizations of $B[1/a]$, are all surjective. In particular,
A.2 holds for every semilocal one-dimensional noetherian domain.
No finiteness of $B$ over $R$ and no excellence is required.
\end{theorem}
\begin{proof}
By Krull--Akizuki in the precise form
\citet[Tag 09IG, Lemma 10.120.18]{A_Stacks}, $B$ is Dedekind and
$\nu:\Spec B\to\Spec R$ is integral and surjective, with finite
fibres and finite residue extensions. In particular $n\geq1$.
When $p=2$, first restrict to a principal nonempty open on which
a fraction-field identity $-1=\sum_i u_i^2$ is defined.
No residue field on this open is orderable. This restriction
does not change the eventual assertion A.2. At a closed point
$x$ of this open, choose a point $y$ of its normalization above
it. The valuation bound gives $\cd_p\kappa(y)\leq n-1$.
Finite-extension invariance of field cohomological dimension,
with the real exception excluded, gives
$\cd_p\kappa(x)\leq n-1$ as well.

We give the needed sheaf comparison even for an infinite
normalization. Pushforward under an integral morphism is exact
on all abelian \'etale sheaves
\citep[Tag 04C2, Lemma 59.43.5]{A_Stacks} and commutes with
arbitrary base change
\citep[Tag 0EYP, Lemma 59.55.4]{A_Stacks}. For a constructible
$\mathbf F_p$-sheaf $\mathcal F$ on $\Spec R$, the unit is a
monomorphism in
\[
 0\longrightarrow\mathcal F\longrightarrow\nu_*\nu^*\mathcal F
 \longrightarrow\mathcal Q\longrightarrow0.
\]
Indeed its pullback along the surjection $\nu$ is a split
monomorphism by the adjunction identity, and $\nu^*$ is exact and
conservative. Integral base change at the generic point shows
$\mathcal Q_{\overline\eta}=0$.

Any $p$-torsion sheaf $\mathcal Q$ with zero generic stalk on this
curve has $H^q(R,\mathcal Q)=0$ for $q>n-1$. To verify this
without asserting that $\mathcal Q$ is constructible, write it
as a filtered colimit of constructible $\mathbf F_p$-sheaves
$\mathcal Q_i$. For every $i$, finiteness of its generic module
and the vanishing of the generic colimit give a later $j$ for
which the transition $\mathcal Q_i\to\mathcal Q_j$ is zero
at the generic point. Its image is constructible, with support
in a finite set of closed points. Its cohomology above $n-1$
vanishes. Every class in the cohomological colimit therefore
vanishes in those degrees, proving the assertion.

It follows that if $\cd_p(B)\leq n$, then
$H^q(R,\mathcal F)=0$ for $q>n$. Conversely exact integral
pushforward gives $\cd_p(B)\leq\cd_p(R)$. These comparisons
also hold on every principal localization with the stated
residue-field bounds.

Suppose A.2 holds for $B$. Its good affine open has complement
a finite set of closed points of $\Spec B$. Their images in
$\Spec R$ form a finite closed set, since $\nu$ is integral.
Choose a nonzero $a\in R$ vanishing at these images and, when
$p=2$, also enforcing the preceding sum-of-squares restriction.
Then $\Spec B[1/a]$ lies inside the good open. The upper bound
holds on every smaller open of this normalization, by
Theorem~\ref{A:Dedekind-Milnor-criterion}. For every affine
$U\subset\Spec R[1/a]$, its normalization is the corresponding
open of $\Spec B[1/a]$, so the sheaf comparison proves
$\cd_p(U)\leq n$. Proposition~\ref{A:spread} then proves A.2
for $R$.

Conversely, A.2 for $R$ supplies a principal open with
$\cd_p(R[1/a])\leq n$. Exact integral pushforward gives
$\cd_p(B[1/a])\leq n$. The residue criterion makes this upper
bound hereditary on smaller opens of $\Spec B[1/a]$, and generic
spreading proves A.2 for $B$. The same argument proves the stated
equivalence with the simultaneous residue condition. Finally, if
$R$ is semilocal, the finite-fibre assertion makes $B$ semilocal.
Corollary~\ref{A:semilocal-Dedekind} applies, and the established
normalization comparison finishes the proof.
\end{proof}

\begin{remark}[What the reduction does and does not settle]
This criterion settles every semilocal curve case and makes the obstruction
for every noetherian integral curve explicit. In generic dimension two it is
a simultaneous tame-symbol cokernel, rather than only a Picard
quotient. One finite deletion does not always make these cokernels
vanish: Theorem~\ref{A:A2-persistent-arithmetic-counterexample} gives
an excellent arithmetic counterexample already in generic dimension
one. The Nagata--Gabber example has a different outcome: its
eventual equality follows from the principal-completion theorem
proved below. The literature search used
the combinations ``semilocal'', ``\'etale cohomological dimension'',
``fraction field'', ``Dedekind'', and ``Iwasawa''; it located the
standard norm-residue and arithmetic inputs, without establishing
priority for this formulation.
\end{remark}
